% from ega2-1-fr.tex
\oldpage[II]{85}

En vertu de (\hyperref[II.4.5.2-fr]{4.5.2, b)}), on peut se borner à prouver
que pour toute partie finie $Z'$ de $\operatorname{Proj}(S)$ et tout
voisinage ouvert $U$ de $Z'$, il existe un élément homogène $f\in S_+$ tel
que $Z'\subset(\operatorname{Proj}(S))_f\subset U$
(\hyperref[II.2.4.1-fr]{2.4.1}). Or, par définition, l'ensemble fermé $Y$,
complémentaire de $U$ dans $\operatorname{Proj}(S)$, est de la forme
$V_+(\mathfrak{I})$, où $\mathfrak{I}$ est un idéal gradué de $S$, ne
contenant pas $S_+$ (\hyperref[II.2.3.2-fr]{2.3.2})~; d'autre part, les
points de $Z'$ sont par définition des idéaux premiers gradués
$\mathfrak{p}_i$ de $S_+$ ne contenant pas $\mathfrak{I}$
(\hyperref[II.2.3.1-fr]{2.3.1}). Il existe donc un élément
$f\in\mathfrak{I}$ n'appartenant à aucun des $\mathfrak{p}_i$ (Bourbaki,
\emph{Alg. comm.}, chap.~II, \S~1, n\textsuperscript{o}~1, prop.~2), et comme
les $\mathfrak{p}_i$ sont gradués, le raisonnement fait \emph{loc. cit.}
montre qu'on peut même supposer $f$ homogène~; cet élément répond alors à la
question.
\end{proof}

\begin{proposition}[4.5.5]
\label{II.4.5.5-fr}
Supposons que $X$ soit un schéma quasi-compact, ou un préschéma dont l'espace
sous-jacent est noethérien. Les conditions a) à c') de
(\hyperref[II.4.5.2-fr]{4.5.2}) sont aussi équivalentes aux suivantes~:
\begin{enumerate}
  \item[\rm{d)}] Pour tout $\mathcal{O}_X$-Module quasi-cohérent
  $\mathcal{F}$ de type fini, il existe un entier $n_0$ tel que, pour tout
  $n\geq n_0$, $\mathcal{F}(n)$ soit engendré par ses sections au-dessus de
  $X$.
  \item[\rm{d')}] Pour tout $\mathcal{O}_X$-Module quasi-cohérent
  $\mathcal{F}$ de type fini, il existe deux entiers $n>0$, $k>0$ tels que
  $\mathcal{F}$ soit isomorphe à un quotient du $\mathcal{O}_X$-Module
  $\mathcal{L}^{\otimes(-n)}\otimes\mathcal{O}_X^k$.
  \item[\rm{d'')}] La propriété d') a lieu pour tout faisceau d'idéaux
  quasi-cohérent de type fini de $\mathcal{O}_X$.
\end{enumerate}
\end{proposition}

\begin{proof}
Comme $X$ est quasi-compact, si un $\mathcal{O}_X$-Module quasi-cohérent de
type fini $\mathcal{F}$ est tel que $\mathcal{F}(n)$ (qui est de type fini)
est engendré par ses sections au-dessus de $X$, $\mathcal{F}(n)$ est alors
engendré par un nombre fini de ces sections
(\hyperref[0.5.2.3-fr]{0, 5.2.3}), donc d) implique d') et il est clair que
d') implique d''). Comme tout $\mathcal{O}_X$-Module quasi-cohérent
$\mathcal{G}$ est limite inductive de ses sous-$\mathcal{O}_X$-Modules de
type fini (\hyperref[I.9.4.9-fr]{I, 9.4.9}), pour vérifier la condition c')
de (\hyperref[II.4.5.2-fr]{4.5.2}), il suffit de le faire pour un faisceau
quasi-cohérent d'idéaux de $\mathcal{O}_X$ qui est de type fini, et d'')
entraîne donc c'). Reste à voir que si $\mathcal{L}$ est ample la propriété
d) est vérifiée. Considérons un recouvrement fini de $X$ par des
$X_{f_i}$ ($f_i\in S_{n_i}$) qu'on peut supposer affines~; en remplaçant les
$f_i$ par des puissances convenables (ce qui ne change pas les $X_{f_i}$),
on peut supposer tous les $n_i$ égaux à un même entier $m$. Le faisceau
$\mathcal{F}|X_{f_i}$, étant de type fini par hypothèse, est engendré par un
nombre fini de ses sections $h_{ij}$ au-dessus de $X_{f_i}$
(\hyperref[I.1.3.13-fr]{I, 1.3.13})~; il existe donc un entier $k_0$ tel que
la section $h_{ij}\otimes f_i^{\otimes k_0}$ se prolonge en une section de
$\mathcal{F}(k_0m)$ au-dessus de $X$ pour tout couple $(i,j)$
(\hyperref[I.9.3.1-fr]{I, 9.3.1}). \emph{A fortiori} les
$h_{ij}\otimes f_i^{\otimes k}$ se prolongent en sections de
$\mathcal{F}(km)$ au-dessus de $X$ pour tout $k\geq k_0$, et pour ces valeurs
de $k$, $\mathcal{F}(km)$ est donc engendré par ses sections au-dessus de
$X$. Pour tout $p$ tel que $0<p<m$, $\mathcal{F}(p)$ est aussi de type fini,
donc il y a un entier $k_p$ tel que
$\mathcal{F}(p)(km)=\mathcal{F}(p+km)$ soit engendré par ses sections
au-dessus de $X$ pour $k\geq k_p$. Prenant $n_0$ plus grand que tous les
$k_pm$, on en conclut que $\mathcal{F}(n)$ est engendré par ses sections
au-dessus de $X$ pour tout $n\geq n_0$, car un tel $n$ s'écrit $n=km+p$ avec
$k\geq k_p$ et $0\leq p<m$.
\end{proof}

\begin{proposition}[4.5.6]
\label{II.4.5.6-fr}
Soient $X$ un schéma quasi-compact, $\mathcal{L}$ un $\mathcal{O}_X$-Module
inversible.
\begin{enumerate}
  \item[\rm{(i)}] Soit $n$ un entier $>0$. Pour que $\mathcal{L}$ soit
  ample, il faut et il suffit que $\mathcal{L}^{\otimes n}$ soit ample.
  \item[\rm{(ii)}] Soit $\mathcal{L}'$ un $\mathcal{O}_X$-Module inversible
  tel que, pour tout $x\in X$, il existe un entier $n>0$

